\documentclass[12pt]{ximera}
\title{Final Exam Review}
\begin{document}
\begin{abstract}
Review for a final on apportionment and game theory: Jefferson, Hamilton, Webster, and Huntington-Hill, the quota rule and all three paradoxes, dominance and saddle points, and mixed strategies.
\end{abstract}
\maketitle

\section*{Final Exam Review}

This review covers apportionment (Hamilton, Jefferson, Adams, Webster, Huntington-Hill,
and the paradoxes) and game theory. For the earlier chapters, see the Exam 1 and Exam 3
reviews.

\begin{problem}
Five states use Jefferson's method to apportion $50$ seats.
\begin{center}
\begin{tabular}{r|rrrrrr}
 & A & B & C & D & E & Total\\\hline
Population & 44 & 159 & 75 & 88 & 34 & 400\\
Standard quota & 5.5 & 19.875 & 9.375 & 11 & 4.25 & 50
\end{tabular}
\end{center}

\begin{prompt}
\textbf{(a)} Standard divisor: $\answer{8}$
\end{prompt}

\begin{prompt}
\textbf{(b)} Which of these modified divisors works for Jefferson's method?
\begin{multipleChoice}
\choice{$8.2$}
\choice{$7.8$}
\choice[correct]{$7.55$}
\choice{$7.3$}
\end{multipleChoice}
\begin{feedback}
Rounding down at $SD=8$ gives only $48$, so the divisor must shrink. Any divisor from
$7.5$ up to about $7.57$ works. At $7.55$ the quotas are $5.83$, $21.06$, $9.93$, $11.66$,
$4.50$, which round down to $5+21+9+11+4=50$. At $7.8$ the total is still only $49$, and at $7.3$ it jumps to $53$.
\end{feedback}
\end{prompt}

\begin{prompt}
\textbf{(c)} Jefferson apportionment: A $\answer{5}$, B $\answer{21}$, C $\answer{9}$,
D $\answer{11}$, E $\answer{4}$
\end{prompt}

\begin{prompt}
\textbf{(d)} Which rule or paradox does this illustrate?
\begin{multipleChoice}
\choice[correct]{A quota-rule violation}
\choice{The Alabama paradox}
\choice{The population paradox}
\choice{The new-states paradox}
\end{multipleChoice}
\begin{feedback}
B's standard quota is $19.875$, so its upper quota is $20$, yet it gets $21$ seats: an
upper-quota violation, typical of Jefferson's bias toward large states.
\end{feedback}
\end{prompt}

\begin{prompt}
\textbf{(e)} Now use Huntington-Hill with the standard divisor. Does it work, and what is
the apportionment?

A $\answer{6}$, B $\answer{20}$, C $\answer{9}$, D $\answer{11}$, E $\answer{4}$
\begin{feedback}
Cutoffs: $\sqrt{30}\approx5.477$ (A rounds up), $\sqrt{380}\approx19.49$ (B up),
$\sqrt{90}\approx9.487$ (C down), D is exactly $11$, and $\sqrt{20}\approx4.47$ (E down).
That totals $50$, so the standard divisor works. The result matches Hamilton's method.
\end{feedback}
\end{prompt}
\end{problem}

\begin{problem}
An ambulance company serves five boroughs with $39$ ambulances, allocated by population
with Hamilton's method.
\begin{center}
\begin{tabular}{r|rrrrrr}
 & A & B & C & D & E & Total\\\hline
Population & 146 & 258 & 130 & 136 & 213 & 883\\
Standard quota & 6.45 & 11.40 & 5.74 & 6.01 & 9.41 & 39
\end{tabular}
\end{center}

\begin{prompt}
\textbf{(a)} Hamilton allocation of $39$: A $\answer{7}$, B $\answer{11}$, C $\answer{6}$,
D $\answer{6}$, E $\answer{9}$
\begin{feedback}
Lower quotas total $37$; the $2$ extras go to C ($0.74$) and A ($0.45$).
\end{feedback}
\end{prompt}

\begin{prompt}
\textbf{(b)} The company buys a $40$th ambulance, and populations are unchanged. Hamilton
allocation of $40$: A $\answer{6}$, B $\answer{12}$, C $\answer{6}$, D $\answer{6}$,
E $\answer{10}$
\begin{feedback}
$SD=\frac{883}{40}=22.075$; quotas $6.61$, $11.69$, $5.89$, $6.16$, $9.65$. The lower
quotas total $37$; the $3$ extras go to C ($0.89$), B ($0.69$), and E ($0.65$).
\end{feedback}
\end{prompt}

\begin{prompt}
\textbf{(c)} What happened?
\begin{multipleChoice}
\choice[correct]{The Alabama paradox: A lost an ambulance when the fleet grew}
\choice{The population paradox}
\choice{The new-states paradox}
\choice{A quota-rule violation}
\end{multipleChoice}
\end{prompt}
\end{problem}

\begin{problem}
The Fairy Kingdom, the Troll Territory, and the Elf Empire apportion $40$ seats by
Hamilton's method.
\begin{center}
\begin{tabular}{r|rrrr}
 & F & T & E & Total\\\hline
Population & 38 & 205 & 99 & 342
\end{tabular}
\end{center}

\begin{prompt}
\textbf{(a)} Standard divisor: $\answer{8.55}$ \qquad Apportionment: F $\answer{4}$,
T $\answer{24}$, E $\answer{12}$
\begin{feedback}
Quotas $4.44$, $23.98$, $11.58$; the lower quotas total $38$, and the $2$ extras go to T
and E.
\end{feedback}
\end{prompt}

\begin{prompt}
\textbf{(b)} The Trolls grow to $212$ and the Elves to $99.01$ (in thousands); the Fairies
stay at $38$. The new standard quotas are $4.355$, $24.297$, $11.348$. New apportionment:
F $\answer{5}$, T $\answer{24}$, E $\answer{11}$
\begin{feedback}
The lower quotas total $39$, and the single extra seat goes to the largest residue: F
($0.355$), just ahead of E ($0.348$).
\end{feedback}
\end{prompt}

\begin{prompt}
\textbf{(c)} What happened?
\begin{multipleChoice}
\choice{The Alabama paradox}
\choice[correct]{The population paradox: the Elves grew but lost a seat to the Fairies, who did not grow}
\choice{The new-states paradox}
\choice{A quota-rule violation}
\end{multipleChoice}
\end{prompt}

\begin{prompt}
\textbf{(d)} With the original populations, does Webster's method at the standard divisor
agree with Hamilton's? 
\begin{multipleChoice}
\choice[correct]{Yes}
\choice{No}
\end{multipleChoice}
\begin{feedback}
Rounding $4.44$, $23.98$, $11.58$ conventionally gives $4+24+12=40$.
\end{feedback}
\end{prompt}
\end{problem}

\begin{problem}
Five states with the populations below (in thousands) use Hamilton's method to apportion
$38$ seats.
\begin{center}
\begin{tabular}{r|rrrrrr}
 & A & B & C & D & E & Total\\\hline
Population & 73.3 & 130 & 64 & 67 & 65 & 399.3\\
Standard quota & 6.98 & 12.37 & 6.09 & 6.38 & 6.19 & 38
\end{tabular}
\end{center}

\begin{prompt}
\textbf{(a)} Standard divisor (two decimals): $\answer[tolerance=0.005]{10.51}$ \qquad
Apportionment: A $\answer{7}$, B $\answer{12}$, C $\answer{6}$, D $\answer{7}$,
E $\answer{6}$
\begin{feedback}
The lower quotas total $36$; the $2$ extras go to A ($0.98$) and D ($0.38$), just ahead of
B ($0.37$).
\end{feedback}
\end{prompt}

\begin{prompt}
\textbf{(b)} A new state F with population $20$ joins. How many seats is a fair addition
for it? $\answer{2}$
\begin{feedback}
$20/10.51\approx1.9$, so $2$ seats, for a total of $40$.
\end{feedback}
\end{prompt}

\begin{prompt}
\textbf{(c)} With $40$ seats, the standard quotas are $6.99$, $12.40$, $6.11$, $6.39$,
$6.20$, $1.91$. New apportionment: A $\answer{7}$, B $\answer{13}$, C $\answer{6}$,
D $\answer{6}$, E $\answer{6}$, F $\answer{2}$
\begin{feedback}
The lower quotas total $37$; the $3$ extras go to A ($0.99$), F ($0.91$), and B ($0.40$),
now just ahead of D ($0.39$).
\end{feedback}
\end{prompt}

\begin{prompt}
\textbf{(d)} What happened?
\begin{multipleChoice}
\choice{The Alabama paradox}
\choice{The population paradox}
\choice[correct]{The new-states paradox: adding F and its fair share of seats moved a seat from D to B}
\choice{A quota-rule violation}
\end{multipleChoice}
\end{prompt}
\end{problem}

\begin{problem}
Player A chooses a row and player B a column; entries are payments from B to A.
\[
\begin{bmatrix} 3 & 5 & 6 & -8\\ 2 & 4 & 7 & 6\\ 5 & 6 & 8 & 5 \end{bmatrix}
\]

\begin{prompt}
\textbf{(a)} Which strategy for A dominates A's first strategy? Row $\answer{3}$
\end{prompt}

\begin{prompt}
\textbf{(b)} After deleting A's first strategy, which strategy for B dominates all of B's
others? Column $\answer{1}$
\begin{feedback}
In the remaining rows, column 1 ($2,5$) is no larger than any other column: $(4,6)$,
$(7,8)$, $(6,5)$. (In the full matrix column 4 is not dominated, because of its $-8$.)
\end{feedback}
\end{prompt}

\begin{prompt}
\textbf{(c)} The saddle point is in row $\answer{3}$, column $\answer{1}$, with value
$\answer{5}$.
\begin{feedback}
The row minimums $-8,2,5$ have maximum $5$; the column maximums $5,6,8,6$ have minimum
$5$.
\end{feedback}
\end{prompt}
\end{problem}

\begin{problem}
You are designing an extension for a botanical garden and must choose mostly
drought-resistant (D), shade-tolerant (S), or pest-resistant (P) trees. The table gives
the percent chance that the trees survive at least $10$ years under each summer weather
pattern.
\begin{center}
\begin{tabular}{l|rrrr}
 & Hot/Humid & Hot/Dry & Cool/Humid & Cool/Dry\\\hline
D & 60 & 90 & 70 & 85\\
S & 70 & 65 & 70 & 80\\
P & 85 & 80 & 75 & 80
\end{tabular}
\end{center}

\begin{prompt}
\textbf{(a)} Planning for Cool/Dry weather, the best choice is:
\begin{multipleChoice}
\choice[correct]{D}
\choice{S}
\choice{P}
\end{multipleChoice}
\end{prompt}

\begin{prompt}
\textbf{(b)} With pest-resistant trees, the ideal weather is:
\begin{multipleChoice}
\choice[correct]{Hot/Humid}
\choice{Hot/Dry}
\choice{Cool/Humid}
\choice{Cool/Dry}
\end{multipleChoice}
\end{prompt}

\begin{prompt}
\textbf{(c)} Which weather column is dominated (worse for ``Weather'' in every row) by
Cool/Humid?
\begin{multipleChoice}
\choice{Hot/Humid}
\choice{Hot/Dry}
\choice[correct]{Cool/Dry}
\end{multipleChoice}
\end{prompt}

\begin{prompt}
\textbf{(d)} Find the saddle point: plant $\answer{P}$, with guaranteed survival chance
$\answer{75}$ percent.
\begin{feedback}
The row minimums are $60$, $65$, $75$, with maximum $75$ (P). The column maximums are $85$,
$90$, $75$, $85$, with minimum $75$ (Cool/Humid). They meet at P, Cool/Humid.
\end{feedback}
\end{prompt}
\end{problem}

\begin{problem}
Aiden and Bailey each have a red and a blue ball and reveal one on the count of three. If
both show blue, Aiden wins $3$ points; if both show red, nobody scores; if Aiden shows red
and Bailey blue, Bailey wins $2$; if Aiden shows blue and Bailey red, Bailey wins $1$.

\begin{prompt}
\textbf{(a)} With rows and columns ordered red, blue, give Aiden's payoff matrix.

$\begin{bmatrix} \answer{0} & \answer{-2}\\ \answer{-1} & \answer{3} \end{bmatrix}$
\end{prompt}

\begin{prompt}
\textbf{(b)} Is the game strictly determined?
\begin{multipleChoice}
\choice{Yes}
\choice[correct]{No}
\end{multipleChoice}
\begin{feedback}
The row minimums $-2,-1$ have maximum $-1$; the column maximums $0,3$ have minimum $0$. No
saddle point.
\end{feedback}
\end{prompt}

\begin{prompt}
\textbf{(c)} At equilibrium, Aiden shows red with probability $\answer{2/3}$, and Bailey
shows red with probability $\answer{5/6}$.
\begin{feedback}
Aiden: $0p-1(1-p)=-2p+3(1-p)$ gives $6p=4$. Bailey: $0q-2(1-q)=-q+3(1-q)$ gives $6q=5$.
\end{feedback}
\end{prompt}

\begin{prompt}
\textbf{(d)} Value of the game: $\answer{-1/3}$
\begin{feedback}
$p-1=\frac23-1=-\frac13$: the game favors Bailey slightly.
\end{feedback}
\end{prompt}
\end{problem}

\begin{problem}
You must decide what to wear on a December day: a full bundle (coat, hat, scarf, mittens),
just a jacket, or no outerwear. Your payoffs depend on the weather:
\[
\begin{array}{r|rrrr}
 & \text{sun} & \text{snow} & \text{rain} & \text{wind}\\\hline
\text{Full bundle} & -1 & 2 & 3 & 2\\
\text{Just a jacket} & 1 & 1 & 1 & 1\\
\text{No outerwear} & 0 & -1 & -3 & -2
\end{array}
\]

\begin{prompt}
\textbf{(a)} Which strategy is dominated?
\begin{multipleChoice}
\choice{Full bundle}
\choice{Just a jacket}
\choice[correct]{No outerwear}
\end{multipleChoice}
\begin{feedback}
Just a jacket beats no outerwear under every kind of weather.
\end{feedback}
\end{prompt}

\begin{prompt}
\textbf{(b)} The forecast is a $50\%$ chance of sun and $50\%$ chance of wind, and you
flip a coin between the full bundle and the jacket. What is your expected payoff?
$\answer{0.75}$
\begin{feedback}
The game reduces to
\[
\begin{array}{r|rr}
 & \text{sun} & \text{wind}\\\hline
\text{Bundle} & -1 & 2\\
\text{Jacket} & 1 & 1
\end{array}
\]
Bundle: $0.5(-1)+0.5(2)=0.5$. Jacket: $1$. Flipping a coin: $0.5(0.5)+0.5(1)=0.75$. (Just
wearing the jacket would do better.)
\end{feedback}
\end{prompt}
\end{problem}

\end{document}
